\documentclass[11pt,a4paper]{article}

\newcommand{\doctitle}{PRS Analytical Pricing Model (MKM-PR-001)}
\newcommand{\docsubtitle}{Wolfram Oracle Benchmark of the Fair-Spread Implementation}
\newcommand{\docversion}{0.1}
\newcommand{\docdate}{19-October-2026}
\newcommand{\docauthor}{David K Kelly}

% MKM Research Labs — Shared Document Header
% Usage: % MKM Research Labs — Shared Document Header
% Usage: % MKM Research Labs — Shared Document Header
% Usage: \input{../shared/mkm_header}  (before \begin{document})
%
% Required commands before \input:
%   \newcommand{\doctitle}{...}
%   \newcommand{\docsubtitle}{...}
%   \newcommand{\docversion}{1.0}
%   \newcommand{\docdate}{13-February-2026}
%   \newcommand{\docauthor}{David K Kelly}

% ─── Packages ───────────────────────────────────────────────────────────────
\usepackage[utf8]{inputenc}
\usepackage{amsmath,amsfonts,amssymb}
\usepackage{listings}
\usepackage{color,xcolor}
\usepackage{graphicx}
\usepackage{hyperref}
\usepackage[toc,page]{appendix}
\usepackage{tabto}
\usepackage{float}
\usepackage{comment}
\usepackage[en-GB]{datetime2}
\usepackage{setspace}
\usepackage{subcaption}
\usepackage{booktabs}
\usepackage{longtable}
\usepackage{adjustbox}
\usepackage{enumitem}
\usepackage{fancyhdr}
\usepackage{lastpage}

% ─── Page geometry ──────────────────────────────────────────────────────────
\onehalfspacing
\usepackage{geometry}
\geometry{a4paper, total={170mm,247mm}, left=20mm, top=25mm, bottom=25mm}

% ─── Headers and footers ───────────────────────────────────────────────────
\pagestyle{fancy}
\fancyhf{}
\fancyhead[L]{\small\textcolor{gray}{MKM Research Labs}}
\fancyhead[R]{\small\textcolor{gray}{\doctitle{} --- v\docversion}}
\fancyfoot[C]{\small\textcolor{gray}{CONFIDENTIAL}}
\fancyfoot[R]{\small\textcolor{gray}{Page \thepage\ of \pageref{LastPage}}}
\renewcommand{\headrulewidth}{0.4pt}
\renewcommand{\footrulewidth}{0.4pt}

% ─── Colours ────────────────────────────────────────────────────────────────
\definecolor{mkmblue}{HTML}{1565C0}
\definecolor{codegreen}{rgb}{0,0.6,0}
\definecolor{codegray}{rgb}{0.5,0.5,0.5}
\definecolor{codepurple}{rgb}{0.58,0,0.82}
\definecolor{backcolour}{rgb}{0.95,0.95,0.92}

% ─── Hyperref setup ─────────────────────────────────────────────────────────
\hypersetup{
    colorlinks=true,
    linkcolor=mkmblue,
    citecolor=mkmblue,
    urlcolor=mkmblue
}

% ─── Code listings ──────────────────────────────────────────────────────────
\lstdefinestyle{mystyle}{
    backgroundcolor=\color{backcolour},
    commentstyle=\color{codegreen},
    keywordstyle=\color{magenta},
    numberstyle=\tiny\color{codegray},
    stringstyle=\color{codepurple},
    basicstyle=\ttfamily\footnotesize,
    breakatwhitespace=false,
    breaklines=true,
    captionpos=b,
    keepspaces=true,
    numbers=left,
    numbersep=5pt,
    showspaces=false,
    showstringspaces=false,
    showtabs=false,
    tabsize=2
}
\lstset{style=mystyle}

% ─── Section formatting ────────────────────────────────────────────────────
\setcounter{tocdepth}{3}
\setcounter{secnumdepth}{3}

% ─── Shared macros ──────────────────────────────────────────────────────────
\newcommand{\mkmtitlepage}{%
    \begin{titlepage}
    \begin{center}
    \vspace*{1cm}
    \begin{figure}[H]\centering
        \includegraphics[width=0.4\textwidth]{../shared/Logos/MKM_logo.png}
    \end{figure}
    \vspace{1.5cm}
    {\Huge\bfseries\textcolor{mkmblue}{\doctitle}\par}
    \vspace{0.5cm}
    {\Large\textcolor{gray}{\docsubtitle}\par}
    \vspace{2cm}
    {\large MKM Research Labs\par}
    \vspace{0.5cm}
    {\large Version \docversion{} --- \docdate\par}
    \vspace{0.5cm}
    {\large \docauthor\par}
    \vspace{2cm}
    \textcolor{red}{\textbf{CONFIDENTIAL --- PROPRIETARY}}
    \end{center}
    \end{titlepage}
}

\newcommand{\mkmlegalpage}{%
    \clearpage
    \thispagestyle{empty}
    \begin{center}
    \large\textbf{Legal Notice}
    \vspace{1em}

    \textbf{PROPRIETARY AND CONFIDENTIAL}
    \end{center}

    \noindent This document, including all algorithms, models, methodology, formulas,
    calculations, and intellectual concepts contained herein, constitutes the exclusive
    intellectual property and confidential trade secrets of MKM Research Labs (``MKM'').

    \vspace{0.5em}
    \noindent Any usage, reproduction, distribution, modification, or disclosure of this
    document or any portion thereof, without the express prior written authorisation from
    MKM Research Labs is strictly prohibited and constitutes an infringement of intellectual
    property rights.

    \vspace{0.5em}
    \noindent The document is provided on an ``as is'' basis. MKM Research Labs makes no
    representations or warranties of any kind, express or implied, regarding its accuracy,
    reliability, or suitability for any particular purpose.

    \vspace{1em}
    \begin{center}
    All rights reserved. \copyright\ 2021--2026 MKM Research Labs.

    \vspace{0.5em}
    Governed by the laws of the United Kingdom.
    \end{center}
    \clearpage
}

\newcommand{\mkmhistorypage}[1]{%
    \section{Document History}
    \begin{table}[ht]
        \centering
        \begin{tabular}{l|l|l|l}
            \textbf{Date} & \textbf{Version} & \textbf{Description} & \textbf{Author} \\
            \hline
            #1
        \end{tabular}
        \label{tab:revision_history}
    \end{table}
    \clearpage
}

% ─── Bank validation protocol section commands ──────────────────────────────
% Use these to ensure every model doc has the required sections.
% They are just styled \section calls — content is provided by each doc.

\newcommand{\mkmexecsummary}{\section{Executive Summary}}
\newcommand{\mkmpurpose}{\section{Model Purpose and Scope}}
\newcommand{\mkmframework}{\section{Mathematical Framework}}
\newcommand{\mkminputs}{\section{Input Parameters and Calibration}}
\newcommand{\mkmimplementation}{\section{Implementation Details}}
\newcommand{\mkmsensitivity}{\section{Sensitivity Analysis}}
\newcommand{\mkmvalidation}{\section{Validation and Backtesting}}
\newcommand{\mkmtestresults}{\subsection{Automated Test Results}\input{test_results}}
\newcommand{\mkmlimitations}{\section{Model Limitations and Known Weaknesses}}

\newcommand{\mkmgovernance}{\section{Governance Validation Questions}
\noindent The following nine questions are mandated by Chapter~5 of the
\emph{Model Risk Governance for Vendors} handbook and must be explicitly
addressed for every model in the inventory. Response status is tracked in the
Model Governance Dashboard (MRC portal). \par\medskip}

\newcommand{\mkmvqitem}[3]{\subsection{Q#1 --- #2}
\noindent\textit{``#3''}\par\medskip\noindent}

\newcommand{\mkmchangelog}{\section{Change History}}
  (before \begin{document})
%
% Required commands before \input:
%   \newcommand{\doctitle}{...}
%   \newcommand{\docsubtitle}{...}
%   \newcommand{\docversion}{1.0}
%   \newcommand{\docdate}{13-February-2026}
%   \newcommand{\docauthor}{David K Kelly}

% ─── Packages ───────────────────────────────────────────────────────────────
\usepackage[utf8]{inputenc}
\usepackage{amsmath,amsfonts,amssymb}
\usepackage{listings}
\usepackage{color,xcolor}
\usepackage{graphicx}
\usepackage{hyperref}
\usepackage[toc,page]{appendix}
\usepackage{tabto}
\usepackage{float}
\usepackage{comment}
\usepackage[en-GB]{datetime2}
\usepackage{setspace}
\usepackage{subcaption}
\usepackage{booktabs}
\usepackage{longtable}
\usepackage{adjustbox}
\usepackage{enumitem}
\usepackage{fancyhdr}
\usepackage{lastpage}

% ─── Page geometry ──────────────────────────────────────────────────────────
\onehalfspacing
\usepackage{geometry}
\geometry{a4paper, total={170mm,247mm}, left=20mm, top=25mm, bottom=25mm}

% ─── Headers and footers ───────────────────────────────────────────────────
\pagestyle{fancy}
\fancyhf{}
\fancyhead[L]{\small\textcolor{gray}{MKM Research Labs}}
\fancyhead[R]{\small\textcolor{gray}{\doctitle{} --- v\docversion}}
\fancyfoot[C]{\small\textcolor{gray}{CONFIDENTIAL}}
\fancyfoot[R]{\small\textcolor{gray}{Page \thepage\ of \pageref{LastPage}}}
\renewcommand{\headrulewidth}{0.4pt}
\renewcommand{\footrulewidth}{0.4pt}

% ─── Colours ────────────────────────────────────────────────────────────────
\definecolor{mkmblue}{HTML}{1565C0}
\definecolor{codegreen}{rgb}{0,0.6,0}
\definecolor{codegray}{rgb}{0.5,0.5,0.5}
\definecolor{codepurple}{rgb}{0.58,0,0.82}
\definecolor{backcolour}{rgb}{0.95,0.95,0.92}

% ─── Hyperref setup ─────────────────────────────────────────────────────────
\hypersetup{
    colorlinks=true,
    linkcolor=mkmblue,
    citecolor=mkmblue,
    urlcolor=mkmblue
}

% ─── Code listings ──────────────────────────────────────────────────────────
\lstdefinestyle{mystyle}{
    backgroundcolor=\color{backcolour},
    commentstyle=\color{codegreen},
    keywordstyle=\color{magenta},
    numberstyle=\tiny\color{codegray},
    stringstyle=\color{codepurple},
    basicstyle=\ttfamily\footnotesize,
    breakatwhitespace=false,
    breaklines=true,
    captionpos=b,
    keepspaces=true,
    numbers=left,
    numbersep=5pt,
    showspaces=false,
    showstringspaces=false,
    showtabs=false,
    tabsize=2
}
\lstset{style=mystyle}

% ─── Section formatting ────────────────────────────────────────────────────
\setcounter{tocdepth}{3}
\setcounter{secnumdepth}{3}

% ─── Shared macros ──────────────────────────────────────────────────────────
\newcommand{\mkmtitlepage}{%
    \begin{titlepage}
    \begin{center}
    \vspace*{1cm}
    \begin{figure}[H]\centering
        \includegraphics[width=0.4\textwidth]{../shared/Logos/MKM_logo.png}
    \end{figure}
    \vspace{1.5cm}
    {\Huge\bfseries\textcolor{mkmblue}{\doctitle}\par}
    \vspace{0.5cm}
    {\Large\textcolor{gray}{\docsubtitle}\par}
    \vspace{2cm}
    {\large MKM Research Labs\par}
    \vspace{0.5cm}
    {\large Version \docversion{} --- \docdate\par}
    \vspace{0.5cm}
    {\large \docauthor\par}
    \vspace{2cm}
    \textcolor{red}{\textbf{CONFIDENTIAL --- PROPRIETARY}}
    \end{center}
    \end{titlepage}
}

\newcommand{\mkmlegalpage}{%
    \clearpage
    \thispagestyle{empty}
    \begin{center}
    \large\textbf{Legal Notice}
    \vspace{1em}

    \textbf{PROPRIETARY AND CONFIDENTIAL}
    \end{center}

    \noindent This document, including all algorithms, models, methodology, formulas,
    calculations, and intellectual concepts contained herein, constitutes the exclusive
    intellectual property and confidential trade secrets of MKM Research Labs (``MKM'').

    \vspace{0.5em}
    \noindent Any usage, reproduction, distribution, modification, or disclosure of this
    document or any portion thereof, without the express prior written authorisation from
    MKM Research Labs is strictly prohibited and constitutes an infringement of intellectual
    property rights.

    \vspace{0.5em}
    \noindent The document is provided on an ``as is'' basis. MKM Research Labs makes no
    representations or warranties of any kind, express or implied, regarding its accuracy,
    reliability, or suitability for any particular purpose.

    \vspace{1em}
    \begin{center}
    All rights reserved. \copyright\ 2021--2026 MKM Research Labs.

    \vspace{0.5em}
    Governed by the laws of the United Kingdom.
    \end{center}
    \clearpage
}

\newcommand{\mkmhistorypage}[1]{%
    \section{Document History}
    \begin{table}[ht]
        \centering
        \begin{tabular}{l|l|l|l}
            \textbf{Date} & \textbf{Version} & \textbf{Description} & \textbf{Author} \\
            \hline
            #1
        \end{tabular}
        \label{tab:revision_history}
    \end{table}
    \clearpage
}

% ─── Bank validation protocol section commands ──────────────────────────────
% Use these to ensure every model doc has the required sections.
% They are just styled \section calls — content is provided by each doc.

\newcommand{\mkmexecsummary}{\section{Executive Summary}}
\newcommand{\mkmpurpose}{\section{Model Purpose and Scope}}
\newcommand{\mkmframework}{\section{Mathematical Framework}}
\newcommand{\mkminputs}{\section{Input Parameters and Calibration}}
\newcommand{\mkmimplementation}{\section{Implementation Details}}
\newcommand{\mkmsensitivity}{\section{Sensitivity Analysis}}
\newcommand{\mkmvalidation}{\section{Validation and Backtesting}}
\newcommand{\mkmtestresults}{\subsection{Automated Test Results}\input{test_results}}
\newcommand{\mkmlimitations}{\section{Model Limitations and Known Weaknesses}}

\newcommand{\mkmgovernance}{\section{Governance Validation Questions}
\noindent The following nine questions are mandated by Chapter~5 of the
\emph{Model Risk Governance for Vendors} handbook and must be explicitly
addressed for every model in the inventory. Response status is tracked in the
Model Governance Dashboard (MRC portal). \par\medskip}

\newcommand{\mkmvqitem}[3]{\subsection{Q#1 --- #2}
\noindent\textit{``#3''}\par\medskip\noindent}

\newcommand{\mkmchangelog}{\section{Change History}}
  (before \begin{document})
%
% Required commands before \input:
%   \newcommand{\doctitle}{...}
%   \newcommand{\docsubtitle}{...}
%   \newcommand{\docversion}{1.0}
%   \newcommand{\docdate}{13-February-2026}
%   \newcommand{\docauthor}{David K Kelly}

% ─── Packages ───────────────────────────────────────────────────────────────
\usepackage[utf8]{inputenc}
\usepackage{amsmath,amsfonts,amssymb}
\usepackage{listings}
\usepackage{color,xcolor}
\usepackage{graphicx}
\usepackage{hyperref}
\usepackage[toc,page]{appendix}
\usepackage{tabto}
\usepackage{float}
\usepackage{comment}
\usepackage[en-GB]{datetime2}
\usepackage{setspace}
\usepackage{subcaption}
\usepackage{booktabs}
\usepackage{longtable}
\usepackage{adjustbox}
\usepackage{enumitem}
\usepackage{fancyhdr}
\usepackage{lastpage}

% ─── Page geometry ──────────────────────────────────────────────────────────
\onehalfspacing
\usepackage{geometry}
\geometry{a4paper, total={170mm,247mm}, left=20mm, top=25mm, bottom=25mm}

% ─── Headers and footers ───────────────────────────────────────────────────
\pagestyle{fancy}
\fancyhf{}
\fancyhead[L]{\small\textcolor{gray}{MKM Research Labs}}
\fancyhead[R]{\small\textcolor{gray}{\doctitle{} --- v\docversion}}
\fancyfoot[C]{\small\textcolor{gray}{CONFIDENTIAL}}
\fancyfoot[R]{\small\textcolor{gray}{Page \thepage\ of \pageref{LastPage}}}
\renewcommand{\headrulewidth}{0.4pt}
\renewcommand{\footrulewidth}{0.4pt}

% ─── Colours ────────────────────────────────────────────────────────────────
\definecolor{mkmblue}{HTML}{1565C0}
\definecolor{codegreen}{rgb}{0,0.6,0}
\definecolor{codegray}{rgb}{0.5,0.5,0.5}
\definecolor{codepurple}{rgb}{0.58,0,0.82}
\definecolor{backcolour}{rgb}{0.95,0.95,0.92}

% ─── Hyperref setup ─────────────────────────────────────────────────────────
\hypersetup{
    colorlinks=true,
    linkcolor=mkmblue,
    citecolor=mkmblue,
    urlcolor=mkmblue
}

% ─── Code listings ──────────────────────────────────────────────────────────
\lstdefinestyle{mystyle}{
    backgroundcolor=\color{backcolour},
    commentstyle=\color{codegreen},
    keywordstyle=\color{magenta},
    numberstyle=\tiny\color{codegray},
    stringstyle=\color{codepurple},
    basicstyle=\ttfamily\footnotesize,
    breakatwhitespace=false,
    breaklines=true,
    captionpos=b,
    keepspaces=true,
    numbers=left,
    numbersep=5pt,
    showspaces=false,
    showstringspaces=false,
    showtabs=false,
    tabsize=2
}
\lstset{style=mystyle}

% ─── Section formatting ────────────────────────────────────────────────────
\setcounter{tocdepth}{3}
\setcounter{secnumdepth}{3}

% ─── Shared macros ──────────────────────────────────────────────────────────
\newcommand{\mkmtitlepage}{%
    \begin{titlepage}
    \begin{center}
    \vspace*{1cm}
    \begin{figure}[H]\centering
        \includegraphics[width=0.4\textwidth]{../shared/Logos/MKM_logo.png}
    \end{figure}
    \vspace{1.5cm}
    {\Huge\bfseries\textcolor{mkmblue}{\doctitle}\par}
    \vspace{0.5cm}
    {\Large\textcolor{gray}{\docsubtitle}\par}
    \vspace{2cm}
    {\large MKM Research Labs\par}
    \vspace{0.5cm}
    {\large Version \docversion{} --- \docdate\par}
    \vspace{0.5cm}
    {\large \docauthor\par}
    \vspace{2cm}
    \textcolor{red}{\textbf{CONFIDENTIAL --- PROPRIETARY}}
    \end{center}
    \end{titlepage}
}

\newcommand{\mkmlegalpage}{%
    \clearpage
    \thispagestyle{empty}
    \begin{center}
    \large\textbf{Legal Notice}
    \vspace{1em}

    \textbf{PROPRIETARY AND CONFIDENTIAL}
    \end{center}

    \noindent This document, including all algorithms, models, methodology, formulas,
    calculations, and intellectual concepts contained herein, constitutes the exclusive
    intellectual property and confidential trade secrets of MKM Research Labs (``MKM'').

    \vspace{0.5em}
    \noindent Any usage, reproduction, distribution, modification, or disclosure of this
    document or any portion thereof, without the express prior written authorisation from
    MKM Research Labs is strictly prohibited and constitutes an infringement of intellectual
    property rights.

    \vspace{0.5em}
    \noindent The document is provided on an ``as is'' basis. MKM Research Labs makes no
    representations or warranties of any kind, express or implied, regarding its accuracy,
    reliability, or suitability for any particular purpose.

    \vspace{1em}
    \begin{center}
    All rights reserved. \copyright\ 2021--2026 MKM Research Labs.

    \vspace{0.5em}
    Governed by the laws of the United Kingdom.
    \end{center}
    \clearpage
}

\newcommand{\mkmhistorypage}[1]{%
    \section{Document History}
    \begin{table}[ht]
        \centering
        \begin{tabular}{l|l|l|l}
            \textbf{Date} & \textbf{Version} & \textbf{Description} & \textbf{Author} \\
            \hline
            #1
        \end{tabular}
        \label{tab:revision_history}
    \end{table}
    \clearpage
}

% ─── Bank validation protocol section commands ──────────────────────────────
% Use these to ensure every model doc has the required sections.
% They are just styled \section calls — content is provided by each doc.

\newcommand{\mkmexecsummary}{\section{Executive Summary}}
\newcommand{\mkmpurpose}{\section{Model Purpose and Scope}}
\newcommand{\mkmframework}{\section{Mathematical Framework}}
\newcommand{\mkminputs}{\section{Input Parameters and Calibration}}
\newcommand{\mkmimplementation}{\section{Implementation Details}}
\newcommand{\mkmsensitivity}{\section{Sensitivity Analysis}}
\newcommand{\mkmvalidation}{\section{Validation and Backtesting}}
\newcommand{\mkmtestresults}{\subsection{Automated Test Results}\input{test_results}}
\newcommand{\mkmlimitations}{\section{Model Limitations and Known Weaknesses}}

\newcommand{\mkmgovernance}{\section{Governance Validation Questions}
\noindent The following nine questions are mandated by Chapter~5 of the
\emph{Model Risk Governance for Vendors} handbook and must be explicitly
addressed for every model in the inventory. Response status is tracked in the
Model Governance Dashboard (MRC portal). \par\medskip}

\newcommand{\mkmvqitem}[3]{\subsection{Q#1 --- #2}
\noindent\textit{``#3''}\par\medskip\noindent}

\newcommand{\mkmchangelog}{\section{Change History}}

% T1 so the underscores in code identifiers print as underscores.
\usepackage[T1]{fontenc}
\usepackage{lmodern}
% Generated by tools.wolfram_oracle -- do not edit.
\newcommand{\bmGenerated}{2026-10-19}
\newcommand{\bmCodeFile}{\texttt{src/models/hazard/prs\_analytical.py}}
\newcommand{\bmCodeSha}{\texttt{a70f65faf22c369c}}
\newcommand{\bmCodeCommit}{\texttt{7bbb63078a59}}
\newcommand{\bmCodeDirty}{}
\newcommand{\bmNumerics}{numpy 1.26.4, scipy 1.15.1, mpmath 1.4.1}
\newcommand{\bmWolfram}{15.0.0 for Mac OS X ARM (64-bit) (May 26, 2026)}
\newcommand{\bmWolframVia}{ via live kernel session (wolframclient)}
\newcommand{\bmWolframConfirms}{Wolfram and the 50-digit reference agree to $4.0\times10^{-30}$, so the reference is confirmed by a second, independent implementation.}
\newcommand{\bmCases}{160}
\newcommand{\bmWorst}{$2.4\times10^{-15}$}
\newcommand{\bmTol}{$1.0\times10^{-9}$}
\newcommand{\bmFloor}{2}
\newcommand{\bmChecks}{7}
\newcommand{\bmOutcome}{all pass}
\newcommand{\bmDiscLow}{75.5}
\newcommand{\bmDiscHigh}{675}
\newcommand{\bmCurveWorst}{9.3}
\newcommand{\bmCurveOver}{0.1}
\newcommand{\bmBasisOver}{0.01}


\begin{document}

\mkmtitlepage
\mkmlegalpage

\tableofcontents
\clearpage

% =========================================================================
\mkmexecsummary
% =========================================================================

This document benchmarks the PRS Analytical Pricing Model (MKM-PR-001): the
function that turns a gauge's annual trigger probability into the fair spread of
a Physical Risk Swap. It asks two questions, and answers each against references
that share none of the model's numerics.

\begin{enumerate}[leftmargin=*]
\item \textbf{Does the code compute the documented formula?} The model document
  (\texttt{prs\_pricing.pdf}, \S3.1) specifies the spread exactly. Over \bmCases{}
  cases --- probabilities from $10^{-4}$ to the cap, tenors of 1 to 10 years, two
  recoveries, a flat rate and the book's yield curve --- the model agrees with a
  50-digit implementation of that formula to \bmWorst{} (tolerance \bmTol).
  \bmWolframConfirms{} The floor, the zero rule, the cap and the curve
  interpolation behave as documented.
\item \textbf{Does the model document describe the model accurately?} Three of its
  statements are numerical, and the benchmark tests them as written
  (Table~\ref{tab:bm-claims}). Two are not met; the third is a finding.
\end{enumerate}

\paragraph{Findings.} The implementation is correct. The documentation is not:
\begin{enumerate}[leftmargin=*]
\item \textbf{Discretisation is understated (\S8.3).} The quarterly scheme prices
  every spread higher than continuous time by a factor of about
  $1+(\lambda+r)\Delta t/2$, at every tenor. That is \bmDiscLow{} bp at worst for
  probabilities up to 0.20, where the document says the scheme is accurate, and
  \bmDiscHigh{} bp above, where it says ``up to 2--3 bp''. At UK flood
  probabilities (around 1--2\%) the overstatement is about a bp
  (Table~\ref{tab:bm-worked}).
\item \textbf{The curve is no longer flat (\S8.4).} The document says the
  discount curve is flat; the book prices on a term structure. The effect is
  below 1 bp for probabilities under \bmCurveOver, as the document says, but
  reaches \bmCurveWorst{} bp.
\item \textbf{The basis compares two different formulae (\S3.3).} The basis code
  takes the gauge leg as $p\times10^4$ and the property leg from Eq.~7. At the
  same probability the two legs differ, so part of every reported basis is an
  artefact of the formulae: over 1 bp from $p=\bmBasisOver$.
\end{enumerate}

\paragraph{Evidence.} \bmChecks{} checks, outcome: \bmOutcome{}
(Table~\ref{tab:bm-checks}).
\begin{itemize}[leftmargin=*,nosep]
\item Wolfram: \bmWolfram{}\bmWolframVia.
\item Code under test: \bmCodeFile{} at commit \bmCodeCommit{}\bmCodeDirty.
\item SHA-256 of that file: \bmCodeSha{}\ldots
\item Numerics: \bmNumerics.
\item Generated \bmGenerated{} by \texttt{python -m tools.wolfram\_oracle pr001}.
\end{itemize}

% =========================================================================
\mkmpurpose
% =========================================================================

\subsection{What is benchmarked}
\texttt{compute\_prs\_spread} and \texttt{interpolate\_yield\_rate} in
\texttt{src/models/hazard/prs\_analytical.py}: the fair spread every PRS in the
book is priced at, and the curve it is discounted on. The module is loaded from
its file alone, and called exactly as the book calls it.

\subsection{What is not}
The basis waterfall's bookkeeping (counts, distances, inverse-distance weights)
is not a calculation a mathematical oracle can check, and is out of scope; only
the consistency of its two spread legs is examined. The delta engine (MKM-DE-001)
and the XVA overlay (MKM-XV-001) consume this function and are benchmarked
separately. This is a benchmark of the implementation (validation question Q8),
not independent validation (Q9).

% =========================================================================
\mkmframework
% =========================================================================

Following the model document (\S3.1): the annual probability $p$, capped at
0.999, gives the hazard rate $\lambda=-\ln(1-p)$; survival is
$S(t)=e^{-\lambda t}$ and discounting $D(t)=e^{-r(t)\,t}$. With quarterly dates
$t_i=i\,\Delta t$, $\Delta t=\tfrac14$, $n=4T$,
\begin{equation}
A=\sum_{i=1}^{n}\Delta t\,S(t_i)\,D(t_i),\qquad
P=\sum_{i=1}^{n}(1-R)\bigl(S(t_{i-1})-S(t_i)\bigr)\,D\!\left(t_i-\tfrac{\Delta t}{2}\right),
\end{equation}
and the fair spread is $s=10^4\,P/A$ bp, floored at \bmFloor{} bp. The rate $r(t)$
is either flat or the book's curve, linear between whole-year points and flat
beyond them.

\paragraph{Continuous time.} With protection paid when the trigger occurs and
premium paid continuously, both legs integrate the same $S(t)D(t)$:
\begin{equation}
s_{\mathrm{c}}=10^4\,\frac{(1-R)\int_0^T\lambda\,S(t)D(t)\,\mathrm{d}t}{\int_0^T S(t)D(t)\,\mathrm{d}t}
=10^4\,(1-R)\,\lambda,
\end{equation}
whatever the curve and tenor. This is the benchmark the document's \S8.3 refers to.

\paragraph{The quarterly scheme in closed form.} On a flat curve,
$S(t_{i-1})-S(t_i)=(e^{\lambda\Delta t}-1)S(t_i)$ and
$D(t_i-\Delta t/2)=e^{r\Delta t/2}D(t_i)$, so the sums cancel:
\begin{equation}
s=10^4\,(1-R)\,\frac{e^{\lambda\Delta t}-1}{\Delta t}\,e^{r\Delta t/2},
\qquad\frac{s}{s_{\mathrm{c}}}\approx1+\frac{(\lambda+r)\,\Delta t}{2}.
\end{equation}
The premium is paid in arrears with no accrual on a trigger, and the protection
half a period earlier; both raise the spread. The fallback of \S3.1.7 can never
apply: $A\ge\Delta t\,S(T)D(T)>0$ for any $p$ up to the cap.

% =========================================================================
\section{Benchmark Method}
% =========================================================================

\subsection{References}
The equations above are evaluated with \texttt{mpmath} at 50 significant digits,
implemented from the model document rather than the code. Wolfram evaluates the
same sums exactly, with the curve as a \texttt{Piecewise} function, from each
input's exact value at 60-digit precision, returning 30 digits.

\subsection{Cases}
Ten probabilities ($10^{-4}$ to 0.999), four tenors (1, 3, 5, 10 years), two
recoveries (0, 0.4) and two curves (flat 3\%, the book's
\texttt{DEFAULT\_YIELD\_CURVE}): \bmCases{} cases. The book trades $R=0$
(\S3.2); 0.4 checks the recovery term.

% =========================================================================
\mkmvalidation
% =========================================================================

% Generated by tools.wolfram_oracle -- do not edit.
\begin{table}[H]\centering\small
\caption{Benchmark checks}\label{tab:bm-checks}
\begin{adjustbox}{max width=\textwidth}
\begin{tabular}{p{70mm}p{50mm}rl}
\toprule
Check & Result & Tolerance & Outcome \\
\midrule
Fair spread, all cases (Eqs.~1--8) & $2.4\times10^{-15}$ & $1.0\times10^{-9}$ & pass \\
Wolfram agrees with the reference & $4.0\times10^{-30}$ & $1.0\times10^{-9}$ & pass \\
Model agrees with Wolfram & $2.4\times10^{-15}$ & $1.0\times10^{-9}$ & pass \\
A zero probability prices at zero (s5.3.1) & 0 bp, no floor & -- & pass \\
The floor binds below 2 bp (Eq. 8) & p = 0.0001 priced at the floor & -- & pass \\
Probabilities above 0.999 price as 0.999 (s3.1.1) & p = 0.9995 prices as p = 0.999 & -- & pass \\
Yield curve: linear between years, flat beyond & max difference 0.0e+00 over 81 times & -- & pass \\
\bottomrule
\end{tabular}
\end{adjustbox}
\end{table}


\subsection{Fair spread}
% Generated by tools.wolfram_oracle -- do not edit.
{\footnotesize\setlength{\tabcolsep}{4pt}\begin{longtable}{rrrrr}
\caption{Fair spread on the book curve, $R=0$: model against reference. The full grid (both curves, both recoveries) is in \texttt{spreads.csv}}\label{tab:bm-spreads}\\
\toprule
$T$ & $p$ & Reference spread (bp) & Model rel.\ err & Model vs W. \\
\midrule
\endhead
1 & 0.0001 & 1.00444736302943 (floor) & 0 & 0 \\
1 & 0.001 & 10.0501279583779 & $3.5\times10^{-16}$ & $3.3\times10^{-16}$ \\
1 & 0.005 & 50.3767552131246 & $1.2\times10^{-15}$ & $1.2\times10^{-15}$ \\
1 & 0.01 & 101.070945427314 & $8.8\times10^{-16}$ & $8.5\times10^{-16}$ \\
1 & 0.02 & 203.426165376566 & $1.0\times10^{-15}$ & $9.6\times10^{-16}$ \\
1 & 0.05 & 518.499278570660 & $8.1\times10^{-16}$ & $7.4\times10^{-16}$ \\
1 & 0.1 & 1072.28484224664 & $2.7\times10^{-16}$ & $2.4\times10^{-16}$ \\
1 & 0.2 & 2304.91250322506 & $4.4\times10^{-16}$ & $4.5\times10^{-16}$ \\
1 & 0.5 & 7601.46838182698 & $1.1\times10^{-16}$ & $1.3\times10^{-16}$ \\
1 & 0.999 & 185747.399879538 & $6.4\times10^{-17}$ & $4.6\times10^{-17}$ \\
3 & 0.0001 & 1.00546574666302 (floor) & 0 & 0 \\
3 & 0.001 & 10.0603116235579 & $2.0\times10^{-16}$ & $1.9\times10^{-16}$ \\
3 & 0.005 & 50.4276694658987 & $9.1\times10^{-16}$ & $9.3\times10^{-16}$ \\
3 & 0.01 & 101.172762500286 & $6.4\times10^{-16}$ & $6.2\times10^{-16}$ \\
3 & 0.02 & 203.629745311317 & $1.7\times10^{-15}$ & $1.8\times10^{-15}$ \\
3 & 0.05 & 519.007644910989 & $7.1\times10^{-16}$ & $6.1\times10^{-16}$ \\
3 & 0.1 & 1073.29834699770 & $4.7\times10^{-16}$ & $5.1\times10^{-16}$ \\
3 & 0.2 & 2306.91521795683 & $6.0\times10^{-16}$ & $5.3\times10^{-16}$ \\
3 & 0.5 & 7605.94277197570 & $1.6\times10^{-16}$ & $1.8\times10^{-16}$ \\
3 & 0.999 & 185747.565606826 & $6.5\times10^{-17}$ & $5.4\times10^{-17}$ \\
5 & 0.0001 & 1.00511329809846 (floor) & 0 & 0 \\
5 & 0.001 & 10.0567958313271 & $2.2\times10^{-15}$ & $2.3\times10^{-15}$ \\
5 & 0.005 & 50.4102837185495 & $6.0\times10^{-16}$ & $6.1\times10^{-16}$ \\
5 & 0.01 & 101.138474129454 & $1.6\times10^{-15}$ & $1.6\times10^{-15}$ \\
5 & 0.02 & 203.563101221979 & $9.1\times10^{-16}$ & $8.6\times10^{-16}$ \\
5 & 0.05 & 518.855511123588 & $8.3\times10^{-16}$ & $8.6\times10^{-16}$ \\
5 & 0.1 & 1073.04190294206 & $7.9\times10^{-17}$ & $6.9\times10^{-17}$ \\
5 & 0.2 & 2306.58521785897 & $2.7\times10^{-16}$ & $2.0\times10^{-16}$ \\
5 & 0.5 & 7606.04114076090 & $4.6\times10^{-17}$ & $2.2\times10^{-17}$ \\
5 & 0.999 & 185747.565607130 & $1.7\times10^{-16}$ & $1.8\times10^{-16}$ \\
10 & 0.0001 & 1.00509617460275 (floor) & 0 & 0 \\
10 & 0.001 & 10.0566227760028 & $2.2\times10^{-15}$ & $2.1\times10^{-15}$ \\
10 & 0.005 & 50.4093789121315 & $4.7\times10^{-16}$ & $4.7\times10^{-16}$ \\
10 & 0.01 & 101.136569926700 & $9.3\times10^{-16}$ & $9.6\times10^{-16}$ \\
10 & 0.02 & 203.558942620238 & $8.2\times10^{-16}$ & $8.4\times10^{-16}$ \\
10 & 0.05 & 518.843044552095 & $4.5\times10^{-16}$ & $4.9\times10^{-16}$ \\
10 & 0.1 & 1073.01359166405 & $2.1\times10^{-16}$ & $2.1\times10^{-16}$ \\
10 & 0.2 & 2306.53611328221 & $2.7\times10^{-16}$ & $2.7\times10^{-16}$ \\
10 & 0.5 & 7606.04561980318 & $1.6\times10^{-16}$ & $1.1\times10^{-16}$ \\
10 & 0.999 & 185747.565607130 & $1.7\times10^{-16}$ & $1.8\times10^{-16}$ \\
\bottomrule
\end{longtable}}


\subsection{The model document's claims}
% Generated by tools.wolfram_oracle -- do not edit.
\begin{table}[H]\centering\small
\caption{The model document's claims, tested}\label{tab:bm-claims}
\begin{adjustbox}{max width=\textwidth}
\begin{tabular}{p{45mm}p{35mm}p{55mm}l}
\toprule
Claim & Document says & Benchmark finds & Outcome \\
\midrule
s8.3 Quarterly vs continuous time & up to 3 bp, only above $p=0.20$ & 75.5 bp ($p\le0.20$); 674.6 bp above & \textbf{not met} \\
s8.4 Flat vs book curve, 5y & $<1$ bp & 9.3 bp; over 1 bp from $p=0.1$ & \textbf{not met} \\
s3.3 Basis: gauge leg $p\times10^4$ vs Eq.~7 & not stated & up to 2606 bp; over 1 bp from $p=0.01$ & finding \\
\bottomrule
\end{tabular}
\end{adjustbox}
\end{table}

% Generated by tools.wolfram_oracle -- do not edit.
\begin{table}[H]\centering\small
\caption{Worked figures at 5 years, $R=0$, in bp}\label{tab:bm-worked}
\begin{adjustbox}{max width=\textwidth}
\begin{tabular}{rrrrrrrr}
\toprule
$p$ & Quarterly & Continuous & $\Delta$ & Rel. & Flat$-$book & $p\times10^4$ & Basis artefact \\
\midrule
0.0001 & 1.00 & 1.00 & +0.00 & +0.38\% & -0.00 & 1.0 & -0.01 \\
0.001 & 10.04 & 10.01 & +0.04 & +0.39\% & -0.01 & 10.0 & -0.06 \\
0.005 & 50.35 & 50.13 & +0.22 & +0.44\% & -0.07 & 50.0 & -0.41 \\
0.01 & 101.01 & 100.50 & +0.50 & +0.50\% & -0.13 & 100.0 & -1.14 \\
0.02 & 203.30 & 202.03 & +1.27 & +0.63\% & -0.26 & 200.0 & -3.56 \\
0.05 & 518.18 & 512.93 & +5.24 & +1.02\% & -0.68 & 500.0 & -18.86 \\
0.1 & 1071.61 & 1053.61 & +18.01 & +1.71\% & -1.43 & 1000.0 & -73.04 \\
0.2 & 2303.47 & 2231.44 & +72.04 & +3.23\% & -3.11 & 2000.0 & -306.59 \\
0.5 & 7596.72 & 6931.47 & +665.25 & +9.60\% & -9.32 & 5000.0 & -2606.04 \\
\bottomrule
\end{tabular}
\end{adjustbox}
\end{table}


\noindent The relative column of Table~\ref{tab:bm-worked} follows
$(\lambda+r)\Delta t/2$; the book curve's rates sit above 3\%, which is why it
prices slightly higher than the flat curve.

% =========================================================================
\mkmlimitations
% =========================================================================

\begin{enumerate}[leftmargin=*]
\item \textbf{Benchmark, not validation.} The code is shown to compute the
  documented formula; whether a constant hazard rate and quarterly pricing suit
  the product is a separate question.
\item \textbf{Documentation.} \S8.3 and \S8.4 of the model document should be
  corrected to the figures above, and the code's docstring, which says the
  protection leg is discounted at $D(t_i)$, to the midpoint the code and the
  document both use.
\item \textbf{Basis legs.} Whether the gauge leg should be priced with Eq.~7, so
  that the basis measures the property against the gauge rather than one formula
  against another, is for the model owner.
\item \textbf{High probabilities.} Near the cap the spread exceeds $10^5$ bp; the
  formula is computed correctly there, but the number has no economic meaning.
\end{enumerate}

\mkmhistorypage{19-Oct-2026 & 0.1 & First issue & David K Kelly \\}

\end{document}
